\begin{abstract}
Offline speech recognition is now practical on commodity CPUs, but an interactive dictation
system cannot be characterized by one word-error-rate measurement on clean speech. We present an
expanded empirical evaluation of \YazSes, a privacy-first offline voice-input system, across
eight local engine/checkpoint configurations, clean and harder speech, repeated decoding,
multiple OS/ISA environments, interaction-specific latency mechanisms, and real annotated
meetings.

On a controlled Azure CPU host, Parakeet TDT 0.6B v2 produced the lowest clean-speech \wer{}
point estimate in the measured 200-utterance matrix at 2.06\%, while Moonshine/tiny produced the
lowest measured \rtf{} at 0.016. Moving to LibriSpeech \texttt{test-other} increased every
measured engine's \wer{}, but by materially different factors.

Repeated decoding exposed a distinct \texttt{large-v3} failure mode. In a dedicated four-repeat
probe on the same 200 hard-speech utterances, \wer{} ranged from 5.46\% to 6.61\%;
substitutions remained 87, deletions 15, and correct-word hits 3,619, while insertions varied
from 101 to 144. Follow-up experiments show that previous-text conditioning is
checkpoint-dependent and that disabling it removes the observed \texttt{large-v3} runaway
continuation mode in this workload, without establishing a universal average-\wer{} improvement.

The broader evaluation also changes system-level conclusions: synthetic onset padding does not
demonstrate recovery of missed speech; CPU streaming can increase final latency; and the former
diarization clustering threshold produced 75.21\% mean \der{} on the 16-recording AMI test split.
A meeting-specific threshold reduced mean \der{} to 26.71\%. These results argue that offline
interactive speech systems should report reproducibility, tail failures, hardware dependence,
and task-specific metrics alongside conventional aggregate \wer{}.
\end{abstract}
